\subsection{正规空间中闭集的局部有限开覆盖}

定理 3. 设 \((\mathbf{A}_t)_{t\in \mathbf{I}}\) 是闭集 \(\mathbf{Y}\) （正规空间 \(\mathbf{X}\) 中的）的局部有限开覆盖。则存在开覆盖 \((\mathbf{B}_t)_{t\in \mathbf{I}}\) （\(\mathbf{Y}\) 的），使得 \(\overline{\mathbf{B}}_t\subset \mathbf{A}_t\) 对每个 \(i\in I\) 成立。

把指标集 I 良序化（集合论，第 III 章，§ 2，第 3 号，定理 1）。我们将用超限归纳法在 X 中定义开集族 \((\mathbf{B}_t)_{t \in \mathbf{I}}\) ，使得 (i) \(\overline{\mathbf{B}}_t \subset \mathbf{A}_t\) 对每个 \(t \in \mathbf{I}\) 成立；(ii) 对每个 \(t \in \mathbf{I}\) ，由诸 \(\mathbf{B}_\lambda\) （满足 \(\lambda \leqslant t\) ）与诸 \(\mathbf{A}_\lambda\) （满足 \(\lambda > t\) ）组成的族是 Y 的开覆盖。设我们已定义了诸 \(\mathbf{B}_t\) （对 \(t < \gamma\) ），使得 (i) 与 (ii) 对一切 \(t < \gamma\) 满足，我们来证明可定义 \(\mathbf{B}_\gamma\) 使 (i) 与 (ii) 对 \(t = \gamma\) 也满足。先证诸 \(\mathbf{B}_t\) （满足 \(t < \gamma\) 者）与诸 \(\mathbf{A}_t\) （满足 \(t \geqslant \gamma\) 者）构成 Y 的覆盖。由假设，对每个 \(x \in \mathbf{Y}\) 只有有限个指标 \(\lambda \in \mathbf{I}\) 使 \(x \in \mathbf{A}_\lambda\) ，记作 \(\lambda_1 < \lambda_2 < \cdots < \lambda_n\) ；设 \(\lambda_h\) 是诸 \(\lambda_i\) 中满足 \(\lambda_i < \gamma\) 的最大者；若 \(h < n\) ，则 \(x \in \mathbf{A}_{\lambda_n}\) 且 \(\lambda_n \geqslant \gamma\) ；若 \(h = n\) ，则归纳假设表明 \(x\) 属于某个 \(\mathbf{B}_\lambda\) （满足 \(\lambda \leqslant \lambda_n < \gamma\) ），从而我们的断言得证。现在令 \(\mathbf{C} = (\complement\mathbf{Y}) \cup (\bigcup_{t < \gamma} \mathbf{B}_t) \cup (\bigcup_{t > \gamma} \mathbf{A}_t)\) ；C 是开的，且由刚才所证有 \(\complement\mathbf{A}_\gamma \subset \mathbf{C}\) ；由正规空间的公理 ( \(\mathrm{O}_{\mathbf{V}}^{\prime\prime}\) ) ，因此存在开集 V 使 \(\complement\mathbf{A}_\gamma \subset \overline{\mathbf{V}} \subset V \subset C\) 。若令 \(\mathbf{B}_\gamma = \complement\overline{\mathbf{V}}\) ，则 \(\overline{\mathbf{B}}_T \subset \complement V \subset A_T\) 且 \(\mathbf{B}_T \cup C = X\) ，于是诸 \(\mathbf{B}_t\) （满足 \(t \leqslant \gamma\) 者）与诸 \(A_t\) （满足 \(t > \gamma\) 者）覆盖 Y。

注。注意我们只用到覆盖 (A.) 逐点有限这一事实，即 X 的每一点只属于有限个集 A。

定义 2. 设 X 是拓扑空间，f 是定义在 X 上的实值函数。f 的支集，记作 Supp (f)，是 X 中使 \(f(x) = 0\) 对一切 \(x \notin S\) 成立的最小闭集 S。

换言之，Supp (f) 是全体 \(x \in X\) 中满足 \(f(x) \neq 0\) 者所成之集在 X 中的闭包；或者再说一遍，它是全体 \(x \in X\) 中这样的点——x 的每个邻域都含一点 y 使 \(f(y) \neq 0\) ——所成的集。

设 \((f_{t})_{t\in I}\) 是 \(X\) 上的一族有限实值函数，其支集构成局部有限族；则和式 \(\sum_{t\in I}f_t(x)\) 对每个 \(x\in X\) 有定义（因为它只含有限个非零项）。有限实值函数 \(x\to \sum_{t\in I}f_t(x)\) 称为族 \((f_{t})\) 的和，记作 \(\sum_{t\in I}f_{t}\) 。若每个 \(f_{t}\) 都连续，则 \(f = \sum_{i \in I} f_i\) 也连续；因为若 \(x\) 是 \(X\) 的任一点，则有邻域 \(V\) （\(x\) 的）只与 \(f_i\) 的支集中的有限个相交，从而存在有限子集 \(H\) （\(I\) 的）使得 \(f(y) = \sum_{i \in H} f_i(y)\) 对一切 \(y \in V\) 成立。

定义 3. 给定族 \((\mathbf{A}_{t})_{t\in I}\) （由拓扑空间 \(\mathbf{X}\) 的子集组成），实值函数族 \((f_t)_{t\in I}\) （定义在 \(\mathbf{X}\) 上的）称为从属于族 \((\mathbf{A}_t)_{t\in I}\) 的，如果 \(\operatorname{Supp}(f_t)\subset \mathbf{A}_t\) 对每个指标 \(t\in I\) 成立。

所谓 \(\mathbf{X}\) 上的连续单位分解，是指任何族 \((f_{\iota})_{\iota \in \mathbf{I}}\) ——由连续实值函数 \(\geqslant 0\) （\(\mathbf{X}\) 上的）组成——其支集构成局部有限族，且使得 \(\sum_{\iota \in \mathbf{I}} f_{\iota}(x) = \mathrm{i}\) 对一切 \(x \in \mathbf{X}\) 成立。

命题 3. 给定任一局部有限开覆盖 \((\mathbf{A}_t)_{t \in \mathbf{I}}\) （正规空间 \(\mathbf{X}\) 的），存在连续单位分解 \((f_t)_{t \in \mathbf{I}}\) （\(\mathbf{X}\) 上的），从属于覆盖 \((\mathbf{A}_t)_{t \in \mathbf{I}}\) 。

由定理 3，存在开覆盖 \((\mathbf{B}_t)_{t \in \mathbf{I}}\) （\(\mathbf{X}\) 的）使得 \(\overline{\mathbf{B}}_t \subset \mathbf{A}_t\) 对每个 \(t \in \mathbf{I}\) 成立，且显然覆盖 \((\mathbf{B}_t)\) 是局部有限的。由公理 \((\mathrm{O}_{\mathbf{V}}'')\) ，对每个 \(t \in \mathbf{I}\) 存在开集 \(\mathbf{C}_t\) 使得 \(\overline{\mathbf{B}}_t \subset \mathbf{C}_t \subset \overline{\mathbf{C}}_t \subset \mathbf{A}_t\) 。由公理 \((\mathrm{O}_{\mathbf{V}})\) ，对每个 \(t \in \mathbf{I}\) 存在连续映射 \(g_t\) （\(\mathbf{X}\) 到 \([0,1]\) 中的），使得 \(g_t(x) = 1\) 于 \(\overline{\mathbf{B}}_t\) 上，且 \(\overline{g}_t\) 的支集含于 \(\overline{\mathbf{C}}_t\) ，从而含于 \(\mathbf{A}_t\) 。由于 \((\mathbf{B}_t)\) 是 \(\mathbf{X}\) 的覆盖，有 \(\sum_{t \in \mathbf{I}} g_t(x) > 0\) 对每个 \(x \in \mathbf{X}\) 成立；若令

\[
f _ {t} (x) = \frac {g _ {t} (x)}{\sum_ {t \in I} g _ {t} (x)}
\]

对一切 \(t \in I\) 与一切 \(x \in X\) 成立，则诸 \(f_t\) 构成从属于覆盖 \((\mathbf{A}_t)\) 的连续单位分解。

推论. 给定任一局部有限开覆盖 \((\mathbf{A}_t)\) ——闭集 \(\mathbf{F}\) （正规空间 \(\mathbf{X}\) 中的）的——，存在一个族 \((f_t)\) ——由连续实值函数 \(\geqslant 0\) （\(\mathbf{X}\) 上的）组成——从属于覆盖 \((\mathbf{A}_t)_{t \in \mathbf{I}}\) ，且使 \(\sum_{t \in \mathbf{I}} f_t(x) = 1\) 对一切 \(x \in \mathbf{F}\) 成立，\(\sum_{t \in \mathbf{I}} f_t(x) \leqslant 1\) 对一切 \(x \in \mathbf{X}\) 成立。

由 \(\mathbb{CF}\) 与诸 \(A_{t}\) 组成的集族是 \(X\) 的局部有限开覆盖。因此存在从属于这个覆盖的连续单位分解，它由族 \((f_{t})_{t\in I}\) （满足 \(\operatorname{Supp}(f_t)\subset A_t\) ，对每个 \(t\in I\) ）以及一个函数 \(g\) （其支集含于 \(F\) 的余集）组成。族 \((f_{t})\) 显然满足所要求的条件。

